\begin{frame}{Crítica ao emprego do método}
\begin{columns}[T,totalwidth=\textwidth]

\column{0.52\textwidth} % Ajustado de 0.58 para 0.52
\begin{itemize}
    \item Dada a heterogeneidade da aplicação do método Case Survey, busca-se as diretrizes do método em artigos seminais.
    \item O procedimento clássico envolve: selecionar casos já existentes na literatura, criar um esquema de codificação, aplicar múltiplos avaliadores e analisar estatisticamente o conjunto codificado.
\end{itemize}

\column{0.43\textwidth} % Ajustado de 0.50 para 0.43 (Soma total = 0.95)
\centering
\begin{tikzpicture}[node distance=0.35cm]
% Reduzido levemente o minimum width de 4.5cm para 3.8cm para não estourar a margem
\node[draw=steel, fill=soft, rounded corners=2pt, minimum width=3.8cm, minimum height=0.7cm, font=\small] (c1) {Casos publicados};
\node[draw=steel, fill=soft, rounded corners=2pt, minimum width=3.8cm, minimum height=0.7cm, below=of c1, font=\small] (c2) {Esquema de codificação};
\node[draw=steel, fill=soft, rounded corners=2pt, minimum width=3.8cm, minimum height=0.7cm, below=of c2, font=\small] (c3) {Múltiplos avaliadores};
\node[draw=steel, fill=soft, rounded corners=2pt, minimum width=3.8cm, minimum height=0.7cm, below=of c3, font=\small] (c4) {Análise estatística};

\foreach \a/\b in {c1/c2,c2/c3,c3/c4}{
    \draw[-{Latex[length=2mm]}, thick] (\a) -- (\b);
}
\end{tikzpicture}

\end{columns}
\end{frame}